\subsection{Cobertura de servicios, medición y límites BCP}\label{sec:sd8-bcp-cobertura-p07}
La clasificación de Tabla~\ref{tab:sd8-bia} se aplica a cada servicio y a sus dependencias. NAV y AMA son críticos por personas en el agua y custodia; ESC por presencia, retiro, autorización y emergencia; CON por despacho seguro; ALE por alerta y avisos urgentes. BOY es alto para programación/registro de inspección y crítico cuando su información sostiene una decisión de seguridad o custodia. VAR y CTR son altos para planificación/acreditación ordinaria, y críticos para habilitación de maniobra, acceso de emergencia o entrega segura. AMB es medio para indicadores y crítico para respuesta a derrame o entrega que requiera evidencia de seguridad; su trazabilidad de retiro es alta. FAC/ADM son altos para cobro/cierre y expedientes; una validación que habilita custodia conserva clase crítica. AUT hereda la clase de la transacción autenticada y sus autorizaciones; su información institucional no transaccional es baja. No se rebaja una dependencia crítica por denominarla auxiliar. Los servicios compartidos de identidad, diario, autorización y observabilidad usados en una función crítica conservan ese objetivo. \parencite[RT-10.02; p.~22]{sd8-bt}

El catálogo de servicio identifica transacción, población de intentos válidos, principio/fin, error, efecto confirmado, indisponibilidad y fuente de telemetría. La disponibilidad mensual es $100(T_m-T_i)/T_m$, con duración observada $T_m$ e indisponible $T_i$, sin excluir mantenimiento o falla del proveedor por conveniencia. La prueba sintética se complementa con experiencia real, solicitudes fallidas y abandonadas; pantalla o nube disponibles no prueban función terminada. Para un mes de treinta días, los presupuestos son 43,2/216/432/864 minutos en clases crítica/alta/media/baja. Un incidente crítico de cuatro horas consume 240 minutos, más que los 43,2 permitidos: RTO y disponibilidad son obligaciones distintas. La captura manual reduce daño pero no se contabiliza como transacción electrónica disponible. \parencites[arts.~20 y 78; pp.~14--15, 40]{sd8-ba}[1.7; RT-10.01; pp.~5, 22]{sd8-bt}

Las duraciones sin enlace 24 h, sin cobertura de terminal 12 h y de Ilque 8 h corresponden a continuidad electrónica e integridad bajo condiciones distintas. Energía, batería, almacenamiento, permiso vigente y red local son entradas separadas. SD4 conserva una brecha de cobertura de decisiones exteriores y permiso operacional residual durante aislamiento: no se emite permiso nuevo desde copia antigua ni se prolonga vigencia. Un permiso que vence antes de completar la autonomía impide declarar esa función cumplida. Debe probarse corte imprevisto con permisos de distinta antigüedad, bajas, reinicio y reloj no fiable. Radio/papel no cierran la brecha electrónica; se mantiene visible antes de aceptar RT-03.10/13. \parencite[cap.~15, autonomía y sincronización RT-03.10/13, p.~33; emplazamiento RT-06.01, p.~34]{sd8-c07}

\subsection{Mostrador sin enlace y doble falla local}\label{sec:sd8-bcp-mostrador-p07}
BC-01 para el peak de hasta noventa zarpes concentrados en dos mañanas aplica la siguiente secuencia, bajo SRE en recuperación y referente de operaciones en seguridad y coordinación:
\begin{enumerate}
\item Confirmar que el problema es WAN, no pérdida de energía, nodo o autoridad local. Mostrar estado de desconexión y antigüedad; conservar autorización local individual vigente y copia disponible de matrícula, contrato, certificados y antecedentes de tripulación, con acceso por finalidad.
\item Identificar embarcación y persona, contrastar vigencia y completar legajo en no más de noventa segundos. Un antecedente faltante/vencido o autorización local no válida mantiene trámite pendiente; no se fabrica documento ni comparte credencial para evitar bloqueo.
\item Registrar ID de trámite, documentos/versiones, operador y horas de preparación/remisión/respuesta. Enviar por radio u otro canal oficialmente autorizado si está disponible; sin canal, guardar remisión pendiente y comunicarlo. La plataforma prepara y registra, la autoridad concede zarpe. Silencio o legajo preparado no equivalen a permiso.
\item Registrar salida física y regreso con evidencia y sus límites de diez segundos/dos interacciones, distintos de los noventa segundos de legajo. El despacho de combustible y la escuela mantienen sus propias autoridades y permisos. La radio no exige pantalla durante maniobra.
\item Al reconectar, sincronizar automáticamente diario y evidencias con secuencia, huella, autoridad e idempotencia, incluyendo nuevas operaciones concurrentes. Consultar respuesta externa incierta por ID antes de repetir remisión/efecto; Calidad verifica recuento, efecto y tiempos completos dentro de treinta minutos tras restituir enlace.
\end{enumerate}
Como sensibilidad, noventa legajos que consuman noventa segundos cada uno suman 8.100 segundos, 135 minutos de trabajo de preparación. No es un tiempo de espera garantizado ni supone que los noventa lleguen juntos. Operaciones/SRE deben comprobar puestos, cola, mezcla y disponibilidad de autorizaciones sin sustituir esa carga por promedio anual. Las funciones externas indisponibles quedan explícitas según BC-01; un pago, firma, consulta oficial o extracción remota no se confirma desde copia antigua. \parencite[14.1; RT-09.01; decisión 21]{sd8-c07}

En BC-02 se detecta fallo mediante salud y transacción de negocio; SRE conserva incidencia y aísla nodo, Datos verifica corte durable y exclusión del escritor, y SRE promueve el redundante de SD4 con permisos/secretos locales disponibles. Calidad prueba recepción, amarra, salida/regreso, escuela y despacho antes de declarar recuperación. No se abre una segunda autoridad para resolver incertidumbre. Si el redundante también falla o no se acredita aislamiento, se suspende efecto inseguro, se protege almacenamiento y se activa captura alternativa; Synaptix repara/restaura y la marina valida función. Un disco RAID no es respaldo ni garantiza recuperación del diario más reciente. Tiempo de detección/promoción/validación y pérdida admisible local deben medirse por función; el RPO regional no autoriza perder hechos locales aún no sincronizados.

\subsection{Registro manual, contacto y retorno al proceso}\label{sec:sd8-bcp-manual-p07}
Para BC-01--05 el titular técnico es SRE; Datos incorpora/conciliará registros, Seguridad protege acceso y Calidad comprueba. El referente competente de marina dirige función de operaciones, escuela, varadero o administración y requiere suplente con autoridad. Los perfiles O/L, D, Q y S de T-15 son sustituciones requeridas, no personas ya disponibles; suplencias técnicas y funcionales y sus cuadrantes continúan pendientes. Proyecto comunica afectación, modalidad activa, pendientes y próxima comprobación al CLIENTE. El registro por escenario se conserva en \path{componentes/protocolos_continuidad.csv}.

La tolerancia de una actuación sin certeza, autorización o condición de seguridad es cero: no se entrega menor/embarcación ni se inicia salida/izaje/despacho inseguro. Se propone como límite de atención técnica crítica cuatro horas desde interrupción, conservando SLA y presupuesto mensual; no autoriza cuatro horas de maniobra sin registro ni convierte manual en continuidad electrónica. Modo manual sólo se usa mientras haya persona competente, documento/autorización vigente, custodia y conciliación posible. Cada cambio de turno transfiere folios, pendientes y restricciones con firma; si falla cualquiera, el referente suspende la función dependiente. Para servicios altos el objetivo técnico sigue ocho horas; duración segura por función/ola requiere ensayo antes de habilitación y no se presume igual a ese máximo contractual.

Escuela registra lista real, alumno/embarcación/monitor, salida, cambio, regreso y retiro anticipado como hechos distintos. Portería comunica retiro inmediatamente por radio al referente y conserva autorización e identidad del receptor; escuela acusa y actualiza presencia antes de nueva salida/entrega. Salud y contacto de emergencia quedan en acceso restringido a finalidad y personal competente, sin difusión por canal general. Se cotejan al relevo de monitor, con lista local vigente; una matrícula no demuestra asistencia ni mensaje ausente demuestra ubicación. La respuesta humana a discrepancia es inmediata y no espera reconciliación diferida. \parencite[decisiones 6--8; restricción 14]{sd8-c07}

Ante aviso masivo fallido, BC-04 conserva padrón local autorizado y versión, canal, destinatario e intento. Se usan los tres canales cuando estén disponibles; falla de un proveedor activa alternativa permitida y no confirma entrega universal. A no acusantes se llama con registro de número, persona, hora, instrucción, respuesta y nueva acción/escalamiento. Guardia/referente coordina radio y contacto seguro de expuestos; no exige registro mientras se doblan amarras. Último envío en diez minutos, entrega, acuse y actuación son eventos distintos. Al restablecer se incorporan intentos y respuestas con ID, evitando duplicar instrucción o silenciar a quien no respondió. \parencite[RT-16.21; restricciones 1 y 4]{sd8-c07}

Cada hoja/registro alternativo tiene folio único, proceso, objeto/persona, hecho, tiempo, operador, permiso y causa. Se custodia bajo acceso restringido y se inventaría al cambio de turno; Datos vincula imagen/original con evento y validación independiente, nunca registra un segundo cargo por importar un comprobante. Las operaciones de origen no se sobreescriben: correspondencia por ID, secuencia, moneda/período y consulta externa antes de repetir. Calidad contrasta universo de folios, cantidades, autorizaciones y cero pérdidas/duplicados/diferencias no explicadas. Sólo después SRE restituye autoridad normal, referente valida proceso y Proyecto comunica retorno y restricciones remanentes. Se conservan originales y retención probatoria; el papel no se reinstala como registro oficial permanente. El cierre del legado conserva salvaguarda mes 6, suplencia y diario de sección~\ref{sec:sd8-legado-p05}; convivencia no depende de recuperar a la única operadora.

\subsection{Puertas DR, respaldo y calendario de pruebas}\label{sec:sd8-bcp-puertas-p07}
Activo-pasivo preparado conserva una autoridad de escritura y evita conflicto multimaestro, frente al mayor aprovisionamiento y complejidad de activo-activo; su justificación técnica y de capacidad es DR-4 de SD4. Santiago/São Paulo del proveedor aprobado en G-REG es el par ofertado, con distancia orientativa urbana y análisis de amenazas comunes de sección~\ref{sec:sd8-drp}. Esa distancia no acredita inmuebles concretos ni disponibilidad real del entorno: antes de recepción se requiere inventario IaC regional, rutas/zonas, permisos de tratamiento y prueba de evento común, incluyendo control AWS, IAM, DNS, claves y corrupción replicada.

La replicación continua incluye ambos almacenes y objetos/versiones, con latido, corte recuperable y alarmas de DR-01; éxito de SQL aislado no acredita documentos. Transacciones en curso quedan confirmadas, rechazadas o pendientes identificadas; no se reejecuta efecto externo incierto sin consulta. Terminales desconectados mantienen su autoridad autorizada y evidencias hasta volver a sincronizar, sin aceptar credencial antigua como nueva facultad. El retorno DR-04 replica desde autoridad actual hacia primario recuperado, comprueba corte/atraso y objetos, excluye escritor y prueba negocio antes de abrir tráfico. Configuración regional de identidad/claves, fencing y capacidad siguen por verificar; no se salva el RTO omitiendo una puerta.

RT-07.05 exige que el procedimiento sea ejecutable por personal del CLIENTE tras transferencia, además de documentado y automatizado en la mayor medida posible. La consola gestionada debe permitir al referente habilitado iniciar la ejecución autorizada, observar estado y resultado y detener/escalar sin administrar AWS ni introducir comandos técnicos. Synaptix conserva ejecución técnica del servicio y soporte; sólo pedir la conmutación por teléfono no acredita ese requisito. La ruta de acceso y autorización durante fallo de federación, permisos, doble control y ejercicio desde CLIENTE quedan pendientes de arquitectura/prueba. Conmutación completamente automática es deseable de RT-07.08; no se impone como cláusula obligatoria ni se elimina validación de integridad. \parencite[RT-07.05/08; p.~17]{sd8-bt}

Respaldos se cifran en reposo y tránsito conforme a SD4/SD5; inventario identifica algoritmo/configuración, medio y ubicación, clave y custodio independiente del conjunto, autorización de descifrado y resultado. Se prueban recuperación de clave desde entorno aislado y rechazo de borrado/modificación inmutable con credenciales administrativas comprometidas. Tres copias deben cubrir el mismo conjunto y corte; réplica que propaga corrupción no cuenta como sustituto de respaldo. Diario/objetos locales pendientes se protegen localmente antes de purga y no se dan por cubiertos por una copia cloud anterior. Restauración se aísla y aplica revocaciones, eliminaciones autorizadas y retención probatoria antes de servir. La frecuencia diaria/PITR 35 días/doce mensuales/local treinta días no rebaja retención de evidencia del caso.

La retención de fuentes incluye escuela de menores hasta 18 años y cinco más, mínimo diez años; salida/regreso/zarpe cinco; contratos/custodia/saldos seis; residuos/derrame seis; consumo cinco; inspección de fondeo vida del elemento y cinco más; combustible cinco; vídeo seis meses. El inventario de restauración debe separar NAV, AMA, ESC, CON, BOY, VAR, CTR, AMB, FAC, ADM, AUT y ALE, con conjuntos compartidos y su retención dominante, frecuencia, volumen, originales y tiempo completo. Los nueve grupos y tiempos de SD5 no son doce dominios: sigue abierto RT-07.13 y no se reparte volumen arbitrariamente para cerrarlo. \parencite[cap.~15, retención RT-05.10, p.~33]{sd8-c07}

\path{componentes/pruebas_continuidad.csv} enlaza cada prueba con riesgo, ejecutor, verificador, paquete, recurrencia, preparación y salida. Antes de producción E1/E2 se requieren DR real, resiliencia, autonomías/reconexión y seguridad ofensiva independiente; en Operación DR dos veces al año y al menos semestral, resiliencia semestral, restauración representativa mensual e intrusión anual. Una restauración mensual mide selección, transferencia, descifrado, reaplicación, permisos y negocio; objeto legible aislado no prueba recuperación de dominio. Resultado registra versión/carga/tiempos/RPO/corte/integridad/incidencias y corrección/repetición. Los ensayos definidos son futuros.

Se vincula preparación a 1.2.3.2/1.2.4.2, verificación de migración a 1.10.1, integración/continuidad a 1.10.2 y Operación a 1.12.3/4/5 de SD7; SD9/T-13 verifica y T-17/18 conserva recepción/calendario. Pruebas destructivas incompatibles usan ambiente aislado o secuencia con restitución de línea base, no el mismo ambiente simultáneamente. Se elige fecha habilitada anterior al vencimiento mensual/semestral/anual, fuera de congelamientos y catorce eventos, con aviso diez días hábiles; marejada suspende ensayo programado y obliga a repropagar reserva/ventana sin mover hitos. No se inventan fechas de eventos aún no confirmadas. Mantenimiento usa versión compatible, avance y retorno controlados de SD6 sin interrupción; excepción requiere justificación, autorización y medición. El calendario y recursos efectivos continúan por conciliar en SD7. Simulacro anual de incidente y restauración granular son deseables de RT-11.21/RT-07.14, sin declararlos mínimos adicionales por inferencia. \parencites[art.~20; art.~21.3; pp.~14--16]{sd8-ba}[RT-10.05--07; cap.~20; p.~22]{sd8-bt}
